\documentclass[aps,prl,twocolumn,superscriptaddress]{revtex4-2}
\usepackage{amsmath,amssymb,graphicx,xcolor}
\usepackage{amsthm} % for theorem environment used in §IV (three-homes theorem)
\newtheorem{theorem}{Theorem}

\begin{document}

\title{The Fine Structure Constant from Spectral Geometry of the Hopf Fibration}

\author{J.~Loutey}
\affiliation{Independent Researcher, Kent, Washington}

\date{February 25, 2026}

\begin{abstract}
We observe that the fine structure constant $\alpha$ satisfies the cubic
equation $\alpha^3 - K\alpha + 1 = 0$, where the coefficient
$K = \pi(42 + \pi^2/6 - 1/40) = 137.036064$ is constructed from
spectral invariants of the Hopf fibration $S^1 \to S^3 \to S^2$: the
degeneracy-weighted Casimir trace of the $S^2$ base ($B = 42$), the
spectral zeta function of the $S^1$ fiber ($F = \zeta(2) = \pi^2/6$),
and a boundary correction from the $S^3$ total space
($\Delta = 1/40$). The selection principle $B/N = \dim(S^3) = 3$ is
proven algebraically to select $n_{\max} = 3$ uniquely. The resulting
$\alpha^{-1} = 137.035987$ agrees with experiment to
$8.8 \times 10^{-8}$ relative error with zero free parameters. A
combinatorial search over $1.92 \times 10^9$ formulas built from the
same spectral ingredients yields a $p$-value of
$5.2 \times 10^{-9}$, establishing that this precision is very
unlikely to be accidental. The cubic is shown to be the characteristic
polynomial of a traceless $\mathbb{Z}_3$-symmetric circulant matrix,
suggesting democratic coupling among the three bundle components.
However, the combination rule $K = \pi(B + F - \Delta)$ remains
conjectural: this is an empirical formula with structural support, not
a first-principles derivation. An internal Sprint-A structural check
(2026-04-18; see \S\ref{sec:spectral_triple}, \S\ref{sec:triple_selection},
and the residual discussion in the Open Questions) sharpens three
features: the external $\pi$ prefactor is the Hopf-bundle measure
$\mathrm{Vol}(S^2)/4 = \mathrm{Vol}(S^3)/\mathrm{Vol}(S^1)$; the cutoff
$n_{\max} = 3$ is triply selected (the $B/N = \dim(S^3)$ principle, the
unique $(+,+,-)$ sign pattern, and the agreement of two independent
canonical $\Delta$ forms at the same integer); and the $(+,+,-)$ pattern
is shape-compatible with Atiyah--Patodi--Singer eta-invariant boundary
subtraction but is not derivable from any specific
spectral-action functional~\cite{alpha_sprint_a_memo}. The working
hypothesis that emerges is that $\alpha$ is a \emph{projection constant}
between three structurally distinct spectral regimes, rather than a
quantity derivable from a single principle.
\end{abstract}

\maketitle

%======================================================================
\section{Introduction}
%======================================================================

Richard Feynman called the fine structure constant
$\alpha = e^2/(4\pi\epsilon_0\hbar c) \approx 1/137.036$ ``one of the
greatest damn mysteries of physics''~\cite{feynman1985}. Despite a
century of quantum theory, $\alpha$ remains an unexplained input to
the Standard Model. Attempts to derive it
numerically---from~$\pi$,~$e$, prime numbers, or Platonic
solids---have uniformly failed~\cite{eddington1935,barrow2002}.
Eddington's ``proton--electron mass ratio'' argument and Wyler's
symmetric-space formula~\cite{wyler1969} both achieved suggestive
precision but were criticized for lacking a physical mechanism
connecting the mathematical structure to electromagnetic
coupling~\cite{robertson1971}.

The present work is situated within the GeoVac program, which
established that the discrete graph Laplacian for hydrogen is
mathematically equivalent to the Laplace--Beltrami operator on the
unit three-sphere~$S^3$, related to the Schr\"odinger equation via
Fock's 1935 stereographic projection~\cite{paper7}. The electron's
bound-state manifold \textit{is}~$S^3$---a result supported by 18
independent symbolic proofs~\cite{paper7}.

This three-sphere carries a unique fiber bundle structure: the Hopf
fibration~\cite{hopf1931,urbantke2003}
\begin{equation}\label{eq:hopf}
  S^1 \;\longrightarrow\; S^3 \;\longrightarrow\; S^2.
\end{equation}
The base $S^2$ parametrizes photon propagation directions, while the
fiber $S^1$ encodes the $U(1)$ electromagnetic gauge phase. We
identify a formula for $\alpha$ built entirely from the spectral
invariants of these three manifolds, validate it statistically against
a large space of alternative formulas, and identify structural
underpinnings (a $\mathbb{Z}_3$-symmetric circulant matrix) that
explain the cubic form. We also characterize precisely what remains
unestablished.

We state clearly at the outset: \textit{this is an empirical formula
with structural support, not a first-principles derivation.} The
combination rule that assembles the spectral invariants into the
coupling constant $K$ is observed, not derived.

An earlier version of this work~\cite{paper2_old} attempted to extract
$\alpha$ from a ``symplectic impedance'' ratio on the electron
lattice, achieving $6.2 \times 10^{-5}$ relative error with a fitted
helical pitch. The present result eliminates all fitted parameters and
improves precision by a factor of~500.


%======================================================================
\section{The Hopf Bundle as Electron--Photon Geometry}
%======================================================================

Fock's 1935 stereographic projection maps the hydrogen momentum
eigenvalue problem onto the unit three-sphere~\cite{fock1935}. The
eigenvalues of the Laplace--Beltrami operator on $S^3$ are
\begin{equation}\label{eq:s3eig}
  \lambda_n = -(n^2 - 1), \qquad n = 1, 2, 3, \ldots
\end{equation}
with degeneracy $g_n = n^2$, reproducing the hydrogen spectrum upon
projection to flat space via the energy-shell constraint
$p_0^2 = -2E$~\cite{paper7}. This identification is exact: the graph
Laplacian on the discrete lattice and the continuum operator on $S^3$
share the same spectrum, as verified by 18 independent symbolic
proofs~\cite{paper7}.

The three-sphere admits a unique fiber bundle structure, the Hopf
fibration~(\ref{eq:hopf}). Concretely, the map
$\pi: S^3 \to S^2$ sends each point $(z_1, z_2) \in \mathbb{C}^2$
with $|z_1|^2 + |z_2|^2 = 1$ to the ratio $z_1/z_2 \in S^2$. Each
fiber $\pi^{-1}(p)$ is a great circle $S^1 \subset S^3$.

We interpret the three components physically:
\begin{itemize}
\item \textbf{Total space $S^3$:} The electron bound-state manifold,
  as established by the Fock projection.
\item \textbf{Base $S^2$:} The photon momentum sphere, parametrizing
  propagation directions of the electromagnetic field.
\item \textbf{Fiber $S^1$:} The $U(1)$ gauge phase of the
  electromagnetic potential.
\end{itemize}

This identification is natural: the Hopf fibration is the geometric
expression of $U(1)$ gauge symmetry acting on the electron's state
space. Electromagnetic coupling is the ``cost'' of projecting $S^3$
onto~$S^2$---the information lost in the fiber.


\subsection{Spectral-triple setting and published lineage}\label{sec:spectral_triple}

The graph-spectral construction used below places naturally in the
framework of \emph{finite spectral triples on graphs} developed by
Marcolli and van Suijlekom~\cite{marcolli_vs_2014}. In their framework
a \emph{gauge network} is a finite spectral triple whose algebra is a
direct sum of matrix algebras indexed by the vertices of a graph, with
a connection variable assigned to each edge; the spectral action
$\mathrm{Tr}\,f(D/\Lambda)$ on such a spectral triple reproduces Wilson
lattice gauge theory on the underlying graph. The discrete $S^3$ Hopf
graph of Paper~7, with the $U(1)$ and $SU(2)$ gauge extensions
developed elsewhere in the GeoVac corpus, is a specific instance of
this general construction. The Perez-S\'anchez
correction~\cite{perez_sanchez_2024, perez_sanchez_2024_comment}
clarifies that the continuum limit of such gauge networks is
Yang--Mills without Higgs.

The present paper is not a stand-alone conjecture about $\alpha$
derived from an isolated mathematical observation: it is a specific
claim within the published graph-spectral-triple lineage. What remains
novel and GeoVac-specific is the \emph{prediction}---no published
framework of this kind predicts a physical coupling constant. The
\emph{machinery} is established; the \emph{prediction} is this
paper's contribution. This distinction frames the open questions of
\S\ref{sec:statistical}--\S VII: the missing derivation of
$K = \pi(B + F - \Delta)$ is a conjecture \emph{within} a well-developed
spectral-triple framework, not a conjecture about whether such a
framework exists.


%======================================================================
\section{Spectral Invariants of the Bundle}
%======================================================================

Each component of the Hopf bundle carries intrinsic spectral data. We
extract one canonical invariant from each.

\subsection{Base manifold $S^2$: Casimir trace}

The eigenvalues of the Laplacian on $S^2$ are $l(l+1)$ with
degeneracy $2l+1$. These are the Casimir eigenvalues of $SO(3)$. We
define the \textit{degeneracy-weighted Casimir trace} as
\begin{equation}\label{eq:casimir}
  B(n_{\max}) = \sum_{n=1}^{n_{\max}} \sum_{l=0}^{n-1}
  (2l+1)\,l(l+1),
\end{equation}
where the outer sum runs over principal shells (each shell $n$
contains angular momenta $l = 0, \ldots, n-1$). This is the total
angular momentum content of all electron states up to
shell~$n_{\max}$, weighted by the magnetic degeneracy of each
subshell.

At $n_{\max} = 3$, the explicit computation gives
\begin{align}
  B(3) &= \underbrace{0}_{(1,0)}
         + \underbrace{0}_{(2,0)} + \underbrace{6}_{(2,1)}
         + \underbrace{0}_{(3,0)} + \underbrace{6}_{(3,1)}
         + \underbrace{30}_{(3,2)} \notag \\
       &= 42,
\end{align}
where subscripts denote $(n, l)$ pairs.


\subsection{Selection principle}

The total number of states through shell $n_{\max}$ is
$N(n_{\max}) = \sum_{n=1}^{n_{\max}} n^2 = n_{\max}(n_{\max}+1)(2n_{\max}+1)/6$.
We prove that the ratio
\begin{equation}\label{eq:selection}
  \frac{B(n_{\max})}{N(n_{\max})} = \dim(S^3) = 3
\end{equation}
is satisfied uniquely at $n_{\max} = 3$.

The inner sum in Eq.~(\ref{eq:casimir}) has a closed form:
\begin{equation}\label{eq:innersum}
  \sum_{l=0}^{n-1}(2l+1)\,l(l+1) = \frac{n^2(n^2-1)}{2}.
\end{equation}
Therefore $B(m) = \sum_{n=1}^{m} n^2(n^2-1)/2$, which, using the
standard power-sum identities $\sum n^2$ and $\sum n^4$, evaluates to
\begin{equation}\label{eq:Bclosed}
  B(m) = \frac{m(m+1)(2m+1)(m+2)(m-1)}{20}.
\end{equation}
The ratio simplifies by cancellation:
\begin{equation}\label{eq:ratioclosed}
  \frac{B(m)}{N(m)} = \frac{3(m+2)(m-1)}{10}.
\end{equation}
Setting $B/N = 3$ gives $(m+2)(m-1) = 10$, i.e.,
$m^2 + m - 12 = (m+4)(m-3) = 0$,
whose unique positive root is $m = 3$. This completes the algebraic
proof: $n_{\max} = 3$ is the \textit{only} positive integer at which
the selection principle holds.

\begin{table}[h]
\centering
\begin{tabular}{cccc}
\hline\hline
$n_{\max}$ & $B(n_{\max})$ & $N(n_{\max})$ & $B/N$ \\
\hline
1 & 0   & 1  & 0.000 \\
2 & 6   & 5  & 1.200 \\
3 & 42  & 14 & 3.000 \\
4 & 162 & 30 & 5.400 \\
5 & 462 & 55 & 8.400 \\
\hline\hline
\end{tabular}
\caption{The selection principle $B/N = 3$ is satisfied uniquely at
$n_{\max} = 3$, as proven algebraically via
Eq.~(\ref{eq:ratioclosed}).}
\label{tab:selection}
\end{table}

This result connects to the Casimir structure of $SO(4)$. In the
Peter--Weyl decomposition of $L^2(S^3)$, each shell~$n$ carries the
$(j,j)$ representation of $SO(4) \cong (SU(2) \times SU(2))/\mathbb{Z}_2$
with $j = (n-1)/2$. The $SO(3) \subset SO(4)$ Casimir has per-shell
average
\begin{equation}\label{eq:casimir_identity}
  \langle C_3 \rangle_n
  = \frac{1}{n^2}\sum_{l=0}^{n-1}(2l+1)\,l(l+1)
  = \frac{n^2-1}{2} = C_4(n),
\end{equation}
which equals the $SO(4)$ Casimir eigenvalue. The selection principle
$B/N = 3$ is therefore equivalent to requiring that the
degeneracy-weighted average of the $SO(4)$ Casimir over shells
$1 \leq n \leq n_{\max}$ equals $\dim(S^3)$.

The base manifold invariant is
\begin{equation}\label{eq:base}
  \boxed{B = \sum_{n=1}^{3}\sum_{l=0}^{n-1}(2l+1)\,l(l+1) = 42.}
  \vphantom{\Bigg|}
\end{equation}


\subsection{Triple selection at $n_{\max}=3$}\label{sec:triple_selection}

The Sprint-A internal verification~\cite{alpha_sprint_a_memo} makes an
observation that sharpens the selection principle above. The cutoff
$n_{\max} = 3$ is selected not by one independent condition but by
three mutually independent ones, all converging at the same integer:
\begin{enumerate}
  \item The Casimir selection principle
    $B(m)/N(m) = \dim(S^3) = 3$, proven above
    (Eq.~\ref{eq:ratioclosed} and Table~\ref{tab:selection}).
  \item The sign pattern of the combination rule. The eight sign
    assignments $\pi(\pm B \pm F \pm \Delta)$ at $m = 3$ give
    relative errors from $\alpha^{-1}$ spanning seven orders of
    magnitude; the $(+,+,-)$ pattern uniquely hits $\alpha^{-1}$
    within $4.8 \times 10^{-7}$ relative, with the next-best
    $(+,+,+)$ pattern at $1.1 \times 10^{-3}$---a separation of
    $3.5$ orders of magnitude at $m = 3$, and no analogous
    sign-pattern hit at any other integer $m$.
  \item The agreement of two canonical $\Delta$ forms. The boundary
    product $\Delta(m) = 1/(|\lambda_m| \cdot N(m-1))$ used below
    and the Dirac single-level degeneracy form
    $\Delta(m) = 1/g_m^{\mathrm{Dirac}} = 1/[2(m+1)(m+2)]$ identified
    in Phase~4H (see \S\ref{sec:formula}) are in general distinct
    rational functions of $m$; they agree at $m = 3$ (both giving
    $1/40$) and at no other integer cutoff.
\end{enumerate}

These three mechanisms are independent: the first is
representation-theoretic; the second is numerical and uses all three
$K$-ingredients; the third tests agreement of two independently derived
invariants. Their convergence at a single finite integer is stronger
support for the selection of $n_{\max} = 3$ than any one alone. It
does not, however, derive the combination rule
$K = \pi(B + F - \Delta)$ itself, which remains conjectural per
Theorem~\ref{thm:three_homes}.


\subsection{Fiber $S^1$: Spectral zeta function}

The eigenvalues of the Laplacian on $S^1$ (a circle of circumference
$2\pi$) are $k^2$ for $k \in \mathbb{Z}$. The spectral zeta function
is
\begin{equation}
  \zeta_{S^1}(s) = 2\sum_{k=1}^{\infty} k^{-2s}
                 = 2\,\zeta_R(2s),
\end{equation}
where $\zeta_R$ is the Riemann zeta function. At $s = 1$:
\begin{equation}\label{eq:fiber}
  F = \zeta_R(2) = \frac{\pi^2}{6}.
\end{equation}
This is the zeta-regularized spectral invariant of the photon gauge
phase circle. Equivalently, $F$ is a Dirichlet series of the Fock
shell degeneracy,
$F = D_{n^2}(d_{\max}) := \sum_{n\geq 1} n^2\, n^{-d_{\max}} =
\zeta_R(2)$, where $n^2$ is the $n$-th $S^3$ shell degeneracy and
$d_{\max}=4$ is the packing lattice maximum degree (Paper~0),
linking $F$ to both the Fock projection and the packing axioms.


\subsection{Total space $S^3$: Boundary correction}

The total space $S^3$ contributes a finite-size correction from
the interplay of its eigenvalue spectrum~(\ref{eq:s3eig}) with the
cumulative state count. At the selected cutoff $n_{\max} = 3$, the
highest eigenvalue is $\lambda_3 = -(3^2 - 1) = -8$ and the
cumulative state count through the first non-trivial shell is
$N(2) = \sum_{n=1}^{2} n^2 = 1 + 4 = 5$. The boundary term is
\begin{equation}\label{eq:boundary}
  \boxed{\Delta = \frac{1}{|\lambda_3| \times N(2)}
                = \frac{1}{8 \times 5} = \frac{1}{40}.}
                \vphantom{\Bigg|}
\end{equation}
The physical significance of this boundary term is not established.
It may be suggestive of UV/IR mixing on the compact manifold, but
this interpretation is speculative.

\emph{Fock coupling connection.}  The Fock projection weight between
adjacent $n$-shells (Paper~7) gives a squared coupling
$c^2(n,l) = \tfrac{1}{16}[1 - l(l+1)/(n(n+1))]$.  At the highest
angular momentum of the $n_{\max} = 3$ cutoff ($n = 4$, $l = 3$):
$c^2(4,3) = \tfrac{1}{16}[1 - 12/20] = 1/40 = \Delta$.  The boundary
term is the Fock coupling at the angular-momentum edge of the
Paper~2 truncation.

\emph{APS-shape observation.} The minus sign on $\Delta$ in
$K = \pi(B + F - \Delta)$ is the only sign-flipped boundary-type
correction appearing in any standard heat-kernel or index-theory
decomposition on $S^3$~\cite{alpha_sprint_a_memo}. The standard
Connes--Chamseddine bulk expansion
$\mathrm{Tr}\,f(D/\Lambda) \sim \sum_k f_k\,\Lambda^{d-k}\,a_k$
produces an all-positive Seeley--DeWitt sequence on round $S^3$
(positive cutoff moments $f_k$, positive bulk curvature invariants
$a_k$); the only rule in the heat-kernel / index-theory literature
surveyed (Vassilevich, Connes--Chamseddine, van Suijlekom,
Connes--van Suijlekom spectral truncations) that puts a $(-)$ next
to bulk $(+)$ contributions is the Atiyah--Patodi--Singer
eta-invariant boundary subtraction. Paper~2's
$\Delta^{-1} = g_3^{\mathrm{Dirac}} = 40$ is
\emph{shape-compatible} with APS-style boundary subtraction (a
finite Dirac mode count at a cutoff) but is not literally an APS
eta-invariant: the $\eta$-invariant of the Camporesi--Higuchi Dirac
operator on $S^3$ vanishes by symmetry. The APS match is therefore
a shape-level structural hint, not a derivation.

\emph{Supertrace sign rule.}  The sign structure of
$K = \pi(B + F - \Delta)$ has a cleaner reading: in the
Connes--Chamseddine spectral-action framework, the
\emph{supertrace} $\mathrm{Str} = \mathrm{Tr}_{S} - \mathrm{Tr}_{D}/4$
assigns opposite signs to bosonic (scalar Laplace--Beltrami) and
fermionic (Dirac) contributions via the standard $(-1)^{F}$ grading.
$B$ and $F$ are both scalar-sector objects (finite Casimir trace and
Fock-degeneracy Dirichlet series, respectively) and enter with~$(+)$.
$\Delta$ is a spinor-sector object (Dirac mode-count boundary term)
and enters with~$(-)$.  The $(+,+,-)$ sign pattern is thus the
$(-1)^{F}$ grading, not an ad~hoc sign choice.

\emph{Euler--Maclaurin boundary reading.}  The identification
$\Delta^{-1} = g_{3}^{\mathrm{Dirac}} = 40$ (Phase~4H, SM-D) gains a
sharper reading from the Euler--Maclaurin (EM) decomposition of the
Dirac mode-count sum.  For any cutoff $N$,
$\sum_{n=0}^{N} g_{D}(n) = \text{integral} + [g_{D}(0) + g_{D}(N)]/2
+ \text{Bernoulli corrections}$; the upper boundary term is
$g_{D}(N)/2 = 2(N+1)(N+2)$.  At $N = n_{\mathrm{CH}} = 3$ this equals
$2 \cdot 4 \cdot 5 = 40 = \Delta^{-1}$.
The EM formula explains \emph{why} a single-level degeneracy appears
as a structural ingredient of~$K$: it is the boundary correction when
summing the spectral action up to the selection-principle cutoff.

%======================================================================
% Paper 2, Section IV — DROP-IN REWRITE (proposed by Track D5)
% Replaces the current §IV "The Formula" (lines 272-322 of
% papers/conjectures/paper_2_alpha.tex) with a structurally
% upgraded version that reflects the Phase 4B-4I sprint series.
% Conjectural status of K = π(B + F − Δ) is preserved.
%
% Status: drop-in proposal, NOT auto-applied. PI to review per
% CLAUDE.md §13.5 (reframing a conjectural paper requires PI approval).
%======================================================================
\section{The Formula and its Three Homes}\label{sec:formula}
%======================================================================

Combining the three spectral invariants of the previous section, we
define the \textit{coupling constant}:
\begin{equation}\label{eq:K}
  K = \pi\!\left(B + F - \Delta\right)
    = \pi\!\left(42 + \frac{\pi^2}{6} - \frac{1}{40}\right).
\end{equation}
Numerically,
\begin{equation}
  K = \pi\,(42 + 1.644934 - 0.025000)
    = \pi \times 43.619934
    = 137.036064.
\end{equation}

\emph{Structural home of the external $\pi$.} The external $\pi$
prefactor in Eq.~(\ref{eq:K}) has a Hopf-bundle measure
identification~\cite{alpha_sprint_a_memo}:
\begin{equation}\label{eq:pi_hopf_measure}
  \pi \;=\; \frac{\mathrm{Vol}(S^2)}{4}
         \;=\; \frac{\mathrm{Vol}(S^3)}{\mathrm{Vol}(S^1)}
         \;=\; a_0^{2},
\end{equation}
where $a_0 = \sqrt{\pi}$ is the zeroth Seeley--DeWitt coefficient on
unit $S^3$. The three forms are algebraically equivalent via the
Hopf submersion identity
$\mathrm{Vol}(S^3) = \tfrac{1}{2}\,\mathrm{Vol}(S^2) \cdot \mathrm{Vol}(S^1)$.
The reading is ``Hopf-base solid-angle per chiral charge'': the
integration measure of a finite-cutoff trace restricted to Hopf-base
data, normalized by the Dirac chirality count. This identification
categorically separates the external $\pi$ (first-order Hopf-measure
sector) from the internal $\pi^2$ of $F = \zeta_R(2)$ (second-order
Riemann-zeta sector) and from the absence of $\pi$ in $B$ and
$\Delta$ (rational sectors). Paper~18's exchange-constant taxonomy
may therefore gain a new entry corresponding to the Hopf-bundle
topological measure, sitting outside the existing operator-order
$\times$ bundle $\times$ vertex-topology grid.

The physical $\alpha$ satisfies a self-consistency condition
\begin{equation}\label{eq:selfconsistent}
  \frac{1}{\alpha} + \alpha^2 = K,
\end{equation}
which, upon multiplication by $\alpha$, yields the depressed cubic
\begin{equation}\label{eq:cubic}
  \boxed{\alpha^3 - K\alpha + 1 = 0.}
\end{equation}
This cubic has three real roots (discriminant $4K^3 - 27 > 0$). The
physical root is the smallest positive one, corresponding to
perturbative QED ($\alpha \ll 1$).

\begin{table}[h]
\centering
\begin{tabular}{lcc}
\hline\hline
Quantity & Value & Source \\
\hline
$K$ & $137.036064$ & Eq.~(\ref{eq:K}) \\
$\alpha^{-1}$ (formula) & $137.035987$ & Eq.~(\ref{eq:cubic}) \\
$\alpha^{-1}$ (CODATA 2018) & $137.035999\ldots$ & Experiment \\
\hline
Relative error & $8.8 \times 10^{-8}$ & \\
\hline\hline
\end{tabular}
\caption{Comparison of the Hopf bundle formula with the experimental
value.}
\label{tab:comparison}
\end{table}

%----------------------------------------------------------------------
\subsection{The structural question and the verdict of the Phase 4B--4I
program}\label{sec:structural_question}
%----------------------------------------------------------------------

The three ingredients $B = 42$, $F = \pi^2/6$, and $\Delta = 1/40$ enter
Eq.~(\ref{eq:K}) as an additive combination. A natural question is
whether they share a common spectral generator: a single construction
on the geometry of the electron (here, a sector of $S^3$) that
simultaneously produces all three. Over the period 2025--2026 we
carried out an extended structural-decomposition program, documented in
detail in the project's change log, to test this hypothesis across a
wide range of spectral-theoretic mechanisms.

The program ran for eight phases (4B through 4I) and eliminated nine
distinct mechanisms:
Hopf-twist spectral comparison on $S^1 \times S^2$;
packing-$\pi$ hypothesis with $\pi$ forced by Paper~0 axioms;
scalar Hopf graph morphism between $S^3$ and $S^2$;
$S^1$ fiber Fock-weight traces (discrete and continuous);
$S^5$ spectral geometry;
$\zeta$-combination search for $\Delta$;
the Standard-Model running / one-loop vacuum polarization origin for
$\Delta$;
and most recently, the Dirac-on-$S^3$ lift of $B$, $F$, and $\Delta$ as
a common generator (Tier 1, 2026).

The program produced four positive structural identifications, each of
which resolves the internal content of one of the three ingredients but
does not unify them:
\begin{enumerate}
  \item $\kappa \leftrightarrow B$ via the Fock momentum-space weight
    (Phase 4B track~$\alpha$-C): the universal kinetic constant
    $\kappa = -1/16$ of Paper~7 has an exact algebraic relationship to
    the Casimir trace $B$ at the selection cutoff $m=3$.
  \item $F$ as the scalar Fock-degeneracy Dirichlet series at the
    packing exponent (Phase 4F track~$\alpha$-J): with $g_n = n^2$ the
    degeneracy of the $n$-th $S^3$ shell,
    \begin{equation}\label{eq:F_dirichlet}
      F \;=\; D_{n^2}(s)\Big|_{s = d_{\max}}
        \;=\; \sum_{n\geq 1}\,\frac{n^2}{n^{s}}\Bigg|_{s=4}
        \;=\; \zeta_R(2) \;=\; \frac{\pi^2}{6},
    \end{equation}
    where $d_{\max}=4$ is the maximum degree in Paper~0's packing
    axioms.
  \item $\Delta^{-1}$ as a single-level Dirac degeneracy on $S^3$
    (Phase 4H track SM-D): with the Camporesi--Higuchi spectrum
    $|\lambda_n|_{\text{Dirac}} = n + 3/2$ and degeneracy
    $g_n^{\text{Dirac}} = 2(n+1)(n+2)$ on unit $S^3$,
    \begin{equation}\label{eq:Delta_dirac}
      \Delta^{-1} \;=\; g_3^{\text{Dirac}} \;=\; 2 \cdot 4 \cdot 5 \;=\; 40,
    \end{equation}
    with the Paper-2 cutoff value $n=3$ being the unique integer at
    which this identity holds.
  \item $\zeta(3)$ natively in the Dirac sector at $s = d_{\max}$
    (Tier 1 track~D3, 2026): the Dirac analog of
    Eq.~(\ref{eq:F_dirichlet}) evaluates to
    \begin{equation}\label{eq:D_dirac_zeta3}
      D_{g_n^{\text{Dirac}}}^{\text{Fock}}(4)
      \;=\; \sum_{m\geq 1}\frac{2m(m+1)}{m^4}
      \;=\; 2\,\zeta_R(2) + 2\,\zeta_R(3)
      \;=\; \frac{\pi^2}{3} + 2\,\zeta_R(3),
    \end{equation}
    a closed form mixing $F$ with the Apéry constant $\zeta(3)$. We
    return to this in \S\ref{sec:zeta3_aside}.
\end{enumerate}

The combined verdict of Phases 4B--4I is stated in the following
summary.

\begin{theorem}[Three homes, structural mystery]\label{thm:three_homes}
The three ingredients of the combination rule $K = \pi(B + F - \Delta)$
have the following canonical spectral homes on $S^3$:
\begin{itemize}
  \item[$(B)$] Finite truncated Casimir trace on the scalar
    Laplace--Beltrami spectrum at the selection cutoff $m = 3$, with
    degeneracy weight $(2l+1)\,l(l+1)$.
  \item[$(F)$] Infinite Dirichlet series of the scalar Fock-shell
    degeneracy at the packing exponent,
    $F = \sum_{n} n^2 \cdot n^{-d_{\max}} = \zeta_R(2)$.
  \item[$(\Delta)$] Reciprocal of the single-level Dirac degeneracy at
    the selection cutoff, $\Delta^{-1} = g_3^{\text{Dirac}} = 40$.
\end{itemize}
These three objects live in two different spectral sectors of $S^3$ (the
scalar and spinor Laplace--Beltrami sectors); they are three different
object types (finite trace, infinite Dirichlet series, single-level
degeneracy); and no spectral construction on $S^3$ simultaneously
generates all three. The K combination rule is therefore a
\emph{cross-sector sum} with three separate algebraic origins, and its
validity to $8.8 \times 10^{-8}$ is a structural mystery in the sense
that the sum $B + F - \Delta$ collapses to $\alpha^{-1}/\pi$ through no
currently known common-generator mechanism.
\end{theorem}

Theorem~\ref{thm:three_homes} is a structural statement about how each
ingredient arises, not a proof of Eq.~(\ref{eq:K}). The combination
rule itself \emph{remains conjectural}. What the Phase 4B--4I program
has achieved is to move the open question from ``how is each piece
derived'' (answered by the four positive identifications above) to
``why does the cross-sector sum close to $\alpha^{-1}/\pi$.''

%----------------------------------------------------------------------
\subsection{The three obstructions to unification}\label{sec:obstructions}
%----------------------------------------------------------------------

The Tier 1 Dirac-on-$S^3$ sprint of 2026 (tracks D1--D5) was the last of
the nine mechanisms tested. Because the Dirac sector was the most
structurally promising (it already contained $\Delta^{-1} = 40$ as a
clean identity), a detailed record of why it also fails to lift $B$ and
$F$ is worth inlining. The arguments are short and closed-form.

\paragraph{Obstruction 1 (to lifting $B$).}
The scalar Laplacian on $S^3$ has a zero mode at $(n=1, l=0)$ whose
Casimir weight $l(l+1)$ vanishes. This gives the closed-form identity
\begin{equation}\label{eq:Bclosed_repeat}
  B(m) \;=\; \frac{m(m-1)(m+1)(m+2)(2m+1)}{20},
\end{equation}
with an explicit factor $(m-1)$ enforced by the zero mode. The Dirac
operator on $S^3$ has no zero mode: $|\lambda_n| = n+3/2 \geq 3/2$, and
the ground-level contribution $g_0^{\text{Weyl}}\cdot|\lambda_0|^k$ is
strictly positive for every $k \geq 0$. Cumulative Dirac/Weyl traces
therefore have closed forms
\begin{equation}
  \sum_{n=0}^{N} g_n^{\text{Dirac}}
  \;=\; \frac{(N+1)(N+2)(2N+3)}{3},
  \qquad
  \sum_{n=0}^{N} g_n^{\text{Weyl}}
  \;=\; \frac{(N+1)(N+2)(2N+3)}{6},
\end{equation}
neither of which carries the $(m-1)$ factor of $B(m)$. The cleanest
single-point coincidence is at Paper~2's cutoff $m=3$:
$|\lambda_{m-1}|\cdot g_{m-1}^{\text{Weyl}} = (7/2)\cdot 12 = 42 = B(3)$,
which holds because $(m-1)(m+2) = 10$ at $m=3$ and at no other integer
$m$.

Paper 2's $B = 42$ therefore lives on the scalar Laplace--Beltrami
sector of $S^3$ and does not lift to the Dirac spectrum as a universal
identity. The single-point hit at $m=3$ is a consequence of the
elementary identity $(m-1)(m+2) = 10$, not a spectral theorem.

\paragraph{Obstruction 2 (to lifting $F$).}
The Dirac-sector Dirichlet series at $s = d_{\max}=4$ has the closed
form of Eq.~(\ref{eq:D_dirac_zeta3}), which mixes $F = \zeta_R(2)$
additively with $\zeta_R(3)$. The mechanism is that the Dirac degeneracy
factorizes as a mixture of two homogeneity classes,
\begin{equation}
  g_m^{\text{Dirac}} \;=\; 2\,m^2 \;+\; 2\,m,
\end{equation}
and at $s = 4$ the $m^2$ subchannel produces $2\zeta(2) = 2F$ (the
scalar F content) while the $m$ subchannel produces $2\zeta(3)$ (a
channel absent from the scalar degeneracy). By Apéry's theorem and
standard zeta-independence results, $\zeta_R(2)$ and $\zeta_R(3)$ are
$\mathbb{Q}$-linearly independent, so no rational combination of
Dirac-sector Dirichlet series collapses to a rational multiple of $F$.
The only way to isolate $F$ from the Dirac Dirichlet series is to
subtract the $2\zeta(3)$ channel by hand, which is equivalent to
projecting onto the $m^2$ sub-weight, i.e.\ returning to the scalar
sector.

An equivalent restatement in the Camporesi--Higuchi half-integer index:
the Hurwitz-regularized Dirac spectral zeta at $s=4$ evaluates to
\begin{equation}\label{eq:D_dirac_ch}
  D_{\text{dirac}}^{\text{CH}}(4)
  \;=\; \sum_{n\geq 0}\frac{2(n+1)(n+2)}{(n+3/2)^{4}}
  \;=\; \pi^{2} - \frac{\pi^{4}}{12},
\end{equation}
an inseparable mixture of $\pi^{2}$ and $\pi^{4}$ with no rational
multiple of $F$ extractable.

Paper 2's $F = \pi^{2}/6$ therefore lives on the scalar Fock-degeneracy
arithmetic at the packing exponent and does not lift to the Dirac
Dirichlet series.

\paragraph{Obstruction 3 (to unifying $\Delta$ with $B$ or $F$).}
The Orduz $S^{1}$-equivariant decomposition of the Dirac spinor bundle
over the Hopf fibration $S^{1} \to S^{3} \to S^{2}$ (2026, Tier 1 track
D4) gives, at the Paper-2 cutoff $n_{\text{CH}} = 3$, a 40-dimensional
spin-Dirac level with a clean charge-parity factorization
\begin{equation}\label{eq:Delta_hopf}
  \Delta^{-1} \;=\; g_{3}^{\text{Dirac}} \;=\; 40
  \;=\; \underbrace{4 \cdot 5}_{\text{half-integer charges}}
      + \underbrace{5 \cdot 4}_{\text{integer charges}}
  \;=\; \underbrace{20}_{(+)\,\text{chirality}}
      + \underbrace{20}_{(-)\,\text{chirality}}.
\end{equation}
This is structurally cleaner than Paper~2's scalar factorization
$\Delta^{-1} = |\lambda_{3}|\cdot N(2) = 8 \cdot 5$: the Dirac
decomposition is a single polynomial $g_{n}^{\text{Dirac}} = 2(n+1)(n+2)$
at a single level, whereas the scalar factorization is a product of two
distinct cumulative objects. However, the Dirac refinement does not
produce $B$ or $F$, and the scalar ``8'' in Paper~2's factorization is
an eigenvalue of $-\Delta_{\text{LB}}$, not a Dirac/spinor quantity.
The two factorizations of 40 are structurally distinct and
noncombinable.

Together, obstructions 1--3 close the Dirac-on-$S^{3}$ avenue for
unifying the three ingredients of $K$. All supporting computations are
exact sympy symbolic identities; no floating-point evaluation enters
the arguments above. Sources:
\texttt{geovac/dirac\_s3.py},
\texttt{debug/dirac\_d2\_memo.md},
\texttt{debug/dirac\_d3\_memo.md},
\texttt{debug/dirac\_d4\_memo.md},
\texttt{docs/dirac\_s3\_verdict.md}.

%----------------------------------------------------------------------
\subsection{Connection to the framework's rigidity theorems}\label{sec:rigidity_connection}
%----------------------------------------------------------------------

Theorem~\ref{thm:three_homes} is consistent with, and sharpens, three
established results elsewhere in the GeoVac framework:

\begin{itemize}
  \item \textbf{Paper~7 (S$^{3}$ scalar proof chain):} the 18 symbolic
    proofs of Paper~7 validate the scalar Laplace--Beltrami on $S^{3}$
    as the geometric home of the Coulomb spectrum. $B$ and $F$ both
    live on this sector, reinforcing that Paper~7's $S^{3}$ machinery
    is the ingredient-generating geometry for two of the three pieces
    of $K$.
  \item \textbf{Paper~23 (Fock projection rigidity theorem):} the $S^{3}$
    conformal projection is unique to the $-Z/r$ potential because only
    $-Z/r$ produces an $l$-independent spectrum within each $n$-shell.
    Paper~23 deferred the Dirac-sector analog as an explicit open
    problem (Explorer Gap \#1, Tier 1 scoping). Theorem~\ref{thm:three_homes}
    is \emph{compatible} with that rigidity theorem but independent of
    it: $\Delta$ lives on the spinor spectrum as a single-level
    degeneracy, not as a manifestation of a Dirac Fock projection.
    A Dirac analog of the Fock rigidity theorem remains open and is
    deferred to Tier 2 or beyond.
  \item \textbf{Paper~24 (HO rigidity theorem):} the Bargmann--Segal
    lattice is bit-exactly $\pi$-free in rational arithmetic, and
    Paper~24 \S V identifies the Coulomb/HO asymmetry as ``calibration
    $\pi$ is structurally Coulomb-specific.'' The Dirac-on-$S^{3}$
    spectrum of Camporesi--Higuchi is \emph{also} bit-exactly $\pi$-free
    in rational arithmetic (Tier 1 D1 module, 51/51 symbolic tests),
    and the $\pi$ content in Eqs.~(\ref{eq:D_dirac_zeta3})
    and~(\ref{eq:D_dirac_ch}) enters only through the infinite-sum
    analysis, not through the spectrum itself. The Dirac sector on
    $S^{3}$ is therefore a \emph{structural twin} of the
    Bargmann--Segal HO lattice in its rational-arithmetic certification,
    and the Paper~2 $\pi$ content in $K = \pi(B + F - \Delta)$ is the
    Coulomb-specific calibration of Paper~24's taxonomy.
\end{itemize}

%----------------------------------------------------------------------
\subsection{An aside: $\zeta(3)$ in the Dirac sector}\label{sec:zeta3_aside}
%----------------------------------------------------------------------

Eq.~(\ref{eq:D_dirac_zeta3}) is the first appearance of the Apéry
constant $\zeta_R(3)$ as a structural object in the GeoVac framework.
It is produced natively by the weight-$m$ subchannel of the Dirac
degeneracy $g_{m}^{\text{Dirac}} = 2m(m+1)$ at the Paper-0 packing
exponent $d_{\max} = 4$.

We emphasize that $\zeta(3)$ is \emph{not} part of the $K$ combination
rule. It is a \emph{byproduct} of the Dirac sector — a transcendental
that appears when one asks ``what is the Dirac analog of $F$?'' and
receives the answer ``$F$ plus $\zeta(3)$, inseparably.''

Structurally, $\zeta(3)$ has the following features within the
framework:

\begin{itemize}
  \item It is transcendental of unknown arithmetic type (Apéry 1978
    proved irrationality; algebraic independence from $\pi$ is open).
  \item It is not reachable from any scalar Laplace--Beltrami spectral
    construction on odd-dimensional spheres, as established in Phase~4E
    track~$\alpha$-I: round-sphere spectral invariants produce
    integer powers of $\pi$ in volumes, pure rationals in Seeley--DeWitt
    coefficients and heat-kernel expansions, and $\zeta_{R}(\text{odd})
    + \log 2$ in spectral determinants, but $\pi^{2}$ and $\zeta(3)$
    never appear \emph{additively alongside a rational}.
  \item It appears for the first time in the framework in
    Eq.~(\ref{eq:D_dirac_zeta3}) as the weight-$m$ subchannel of a
    first-order spectral operator at the packing exponent.
\end{itemize}

This suggests — and we flag for future work — that the Paper~18
exchange-constant taxonomy (intrinsic / calibration / embedding / flow)
may benefit from a distinction between \emph{even-zeta content} (arising
from second-order Laplace--Beltrami operators through Jacobi-$\theta$
inversion, which systematically lowers the power of $\pi$ by $1/2$ at
each asymptotic order) and \emph{odd-zeta content} (arising from
first-order spectral operators through half-integer Hurwitz
identities). We do not add a tier to Paper~18 here; we simply record
the phenomenology.

$\zeta(3)$ has no role in reproducing $\alpha^{-1}$. Its inclusion in
the present section is for completeness: the Dirac-on-$S^{3}$ sector
carries new content beyond what $K$ uses, and that content is worth
naming.

%----------------------------------------------------------------------
\subsection{Closure of the structural decomposition program}\label{sec:closure}
%----------------------------------------------------------------------

With Theorem~\ref{thm:three_homes} and the three obstructions of
\S\ref{sec:obstructions}, the structural-decomposition program for
$\alpha$ that began in 2025 and ran through Phases 4B--4I reaches its
natural stopping point. The program has achieved:

\begin{enumerate}
  \item A canonical spectral home for each of $B$, $F$, $\Delta$,
    established in closed form.
  \item A complete enumeration of the mechanisms by which these three
    ingredients might share a common spectral generator, with nine
    mechanisms eliminated and no surviving candidate.
  \item Identification of $\zeta(3)$ as a new framework transcendental
    in the Dirac sector, distinct from the Riemann-zeta-at-even-integer
    content of the scalar sector.
  \item The Seeley--DeWitt cancellation theorem (supertrace sprint,
    2026): on unit $S^{3}$, the asymptotic Dirac mode density is
    $g_{n}^{\mathrm{Dirac}} \sim 4n^{2}$ while the scalar density is
    $g_{n}^{\mathrm{scalar}} = n^{2}$, giving a universal ratio
    $a_{k}^{D^{2}} / a_{k}^{\Delta_{\mathrm{LB}}} = 4 = \dim(S)$
    at every Seeley--DeWitt order $k$, where $S$ is the spinor bundle.
    The perturbative Connes--Chamseddine supertrace
    $\mathrm{Str}[f(D^{2}/\Lambda^{2})]$ therefore vanishes identically
    in the asymptotic expansion.  All three ingredients of
    $K/\pi = B + F - \Delta$ are invisible to the perturbative spectral
    action: $B$ is a finite Casimir trace, $F$ is an arithmetic
    Dirichlet series, and $\Delta$ is an Euler--Maclaurin boundary term.
    $K/\pi$ lives \emph{entirely} in the non-perturbative sector of the
    spectral-action supertrace.
\end{enumerate}

The program has \emph{not} achieved a derivation of Eq.~(\ref{eq:K})
from first principles. The agreement of
\begin{equation}
  \pi\!\left(42 + \frac{\pi^{2}}{6} - \frac{1}{40}\right)
\end{equation}
with $\alpha^{-1}$ to $8.8 \times 10^{-8}$ — that is, 7 decimal digits,
which cannot be reproduced in the extensive combinatorial search of
\S\ref{sec:statistical} — remains a structural coincidence of unknown
origin. We have made the mystery sharper (three spectral homes across
two sectors of $S^{3}$), but we have not resolved it.

Accordingly, the result of this paper retains its \emph{conjectural}
status. Future work on the structural origin of $\alpha$ within the
geometric-vacuum framework is, at the time of this writing, on hold.
The next research direction in the project is the spin-full
generalization of the composed qubit encoding (Paper~14), which builds
on the Dirac-on-$S^{3}$ infrastructure developed in the Tier 1 sprint
but is independent of the $\alpha$ program. That direction has a clean
engineering motivation (heavy-element relativistic chemistry) and does
not require the $\alpha$ conjecture to be resolved.

%======================================================================
% End of Paper 2 §IV rewrite (Tier 1 closure, 2026-04-15)
%======================================================================

%======================================================================
\section{Statistical Validation}\label{sec:statistical}
%======================================================================

A formula involving $\pi$, $42$, and $\pi^2/6$ matching $\alpha^{-1}$
to $8.8 \times 10^{-8}$ could in principle be a numerical
coincidence. To test this, we performed a systematic combinatorial
search~\cite{phase2_audit}.

\subsection{Search space}

The search considers formulas of the form $\alpha^{-1} = T(A + B + C)$
or, equivalently, the cubic $\alpha^3 - T(A+B+C)\alpha + 1 = 0$,
where:
\begin{itemize}
\item $A$, $B$, $C$ are drawn from a pool of spectral quantities:
  eigenvalues $|\lambda_n|$, degeneracies $n^2$ and $2l+1$, Casimir
  values $l(l+1)$, cumulative counts $N(n)$, their pairwise products
  up to~50, and the transcendentals $\zeta(2)$, $\zeta(3)$, $\zeta(4)$,
  $\ln 2$, $\pi$, $\pi^2$ (743 unique term values total).
\item $T$ ranges over 14 transcendental prefactors: $\pi$, $2\pi$,
  $\pi^2$, $\pi/2$, $4\pi$, $4\pi^2$, $1/(2\pi)$, $1$, $2$,
  $1/2$, $\pi/4$, $\pi/6$, $1/\pi$, $\sqrt{\pi}$.
\end{itemize}

Both linear ($\alpha^{-1} = T \cdot S$) and cubic
($\alpha^3 - T \cdot S \cdot \alpha + 1 = 0$) interpretations are
tested, giving $1.92 \times 10^9$ total formulas.

\subsection{Results}

\begin{table}[h]
\centering
\begin{tabular}{rrrr}
\hline\hline
Threshold & Hits & Fraction \\
\hline
$10^{-2}$          & 1{,}283{,}592  & $6.7 \times 10^{-4}$ \\
$10^{-3}$          &   125{,}027    & $6.5 \times 10^{-5}$ \\
$10^{-4}$          &    14{,}476    & $7.5 \times 10^{-6}$ \\
$10^{-5}$          &       682      & $3.5 \times 10^{-7}$ \\
$10^{-6}$          &        37      & $1.9 \times 10^{-8}$ \\
$10^{-7}$          &        11      & $5.7 \times 10^{-9}$ \\
$8.8\times 10^{-8}$ &       10      & $5.2 \times 10^{-9}$ \\
\hline\hline
\end{tabular}
\caption{Hit distribution across precision thresholds. The hit
fraction scales approximately linearly with threshold, consistent with
a uniform distribution in $\alpha^{-1}$ space.}
\label{tab:hits}
\end{table}

At the paper's precision of $8.8 \times 10^{-8}$, only 10 formulas
out of $1.92 \times 10^9$ match as well or better, yielding
\begin{equation}
  \boxed{p\text{-value} = 5.2 \times 10^{-9}.}
\end{equation}
The 10 formulas that beat the paper's precision involve combinations
such as $4\pi(5/49 + 14/15 + \pi^2)$ that lack physical
interpretation in the Hopf bundle context. The paper's formula is the
only one among the top matches with a coherent geometric narrative.

The $p$-value establishes that the precision is very unlikely to be
accidental. It does not, by itself, constitute a derivation---but it
does distinguish this formula from generic numerology.


%======================================================================
\section{The Circulant Structure}
%======================================================================

The cubic $\alpha^3 - K\alpha + 1 = 0$ is not an arbitrary polynomial
ansatz. It arises naturally as the characteristic polynomial of a
traceless $3 \times 3$ circulant matrix.

\subsection{Circulant eigenvalue equation}

Consider the traceless circulant
\begin{equation}\label{eq:circulant}
  M = \begin{pmatrix} 0 & b & c \\ c & 0 & b \\ b & c & 0
      \end{pmatrix},
\end{equation}
which has $\mathbb{Z}_3$ symmetry (invariance under cyclic permutation
of rows and columns). Its characteristic polynomial is
\begin{equation}
  \det(M - sI) = -s^3 + 3bc\,s + (b^3 + c^3) = 0.
\end{equation}
Setting $bc = K/3$ and $b^3 + c^3 = -1$ (equivalently,
$\det M = -1$), this becomes
\begin{equation}
  s^3 - Ks + 1 = 0,
\end{equation}
which is \textit{exactly} the paper's cubic with $s = \alpha$. The
eigenvalues of $M$ are the three roots of the cubic.

\subsection{Physical interpretation}

The $\mathbb{Z}_3$ symmetry of the circulant represents
\textit{democratic coupling}: each of the three objects (labeled $1$,
$2$, $3$) interacts equally with the other two. In the Hopf bundle
context, the three objects are the fiber $S^1$, the base $S^2$, and
the total space $S^3$. Each component contributes a spectral
invariant; the circulant structure implies that their coupling respects
the cyclic symmetry of the bundle decomposition.

The $SU(2)$ structure constants $\varepsilon_{ijk}$ exhibit the same
$\mathbb{Z}_3$ cyclic symmetry, producing a circulant with $b = c = 1$
and characteristic polynomial $s^3 - 3s - 2 = 0$. The paper's cubic
has the \textit{same algebraic structure} with coupling strength
$K/3 \approx 45.7$ instead of~1.

\subsection{Structural content and limitations}

The circulant interpretation is a \textit{structural} result: it
explains why the self-consistency equation takes the form of a
depressed cubic with unit constant term. These features follow from
three conditions:
\begin{enumerate}
\item Tracelessness: $\mathrm{tr}(M) = 0$ (eliminates the $s^2$ term).
\item $\mathbb{Z}_3$ symmetry: only two free parameters ($b$, $c$).
\item Unit determinant condition: $\det(M) = b^3 + c^3 = -1$ (fixes
  the constant term).
\end{enumerate}

The condition $\det(M) = -1$ is consistent with the chirality of the
Hopf map (linking number $+1$) and with the fermionic sign convention
($\psi \to -\psi$ under $2\pi$ rotation of the $S^1$ fiber), but
neither identification constitutes a derivation.

The circulant explains \textit{why a cubic}, but does not derive
\textit{why this specific~$K$}. The single free parameter $bc = K/3$
encodes the coupling strength, and its value
$\pi(42 + \pi^2/6 - 1/40)/3$ is not determined by the circulant
structure alone.

\subsection{Forced Hermiticity and the physical phase}

The coupling parameters $b$, $c$ are necessarily complex conjugates.
Setting $s = b + c$ and $p = bc = K/3$, the constraint $b^3 + c^3 = -1$
reduces to $s^3 - Ks + 1 = 0$---the \textit{same cubic} as for
$\alpha$. For real $b$ and $c$, one requires $s^2 \geq 4K/3$, but all
three roots satisfy $s_i^2 < 4K/3 \approx 182.7$ (the largest is
$s_1^2 \approx 137.1$). Therefore~\cite{phase3_audit}
\begin{equation}\label{eq:polar}
  b = \sqrt{K/3}\; e^{i\theta}, \qquad
  c = \sqrt{K/3}\; e^{-i\theta},
\end{equation}
with universal modulus $|b| = |c| = \sqrt{K/3} \approx 6.76$. The
determinant condition fixes the phase:
\begin{equation}\label{eq:phase}
  \cos(3\theta) = -\frac{1}{2(K/3)^{3/2}} \approx -0.00162.
\end{equation}
Three solutions separated by ${\sim}120^\circ$ correspond to the
three roots. The physical root $\alpha$ has $\theta \approx 3\pi/2$
(nearly pure imaginary coupling: $b \approx -i\sqrt{K/3}$).

The circulant $M$ is therefore \textit{Hermitian}, not real symmetric,
consistent with the $U(1)$ gauge structure of the Hopf fiber: the
off-diagonal coupling carries a phase, as expected for a connection on
a principal $U(1)$-bundle.

\subsection{Self-referential property}

The fact that $b + c$ satisfies the same cubic as $\alpha$ is
structural, not coincidental. For any circulant of the
form~(\ref{eq:circulant}), the $k = 0$ discrete Fourier transform (DFT)
mode is $\lambda_0 = b + c$, which is by construction an eigenvalue
of~$M$ and therefore satisfies its characteristic polynomial. All three
eigenvalues of $M$ are the three cubic roots, regardless of which root
$b + c$ equals~\cite{phase3_audit}.

The circulant decomposition does not factor $K$ into independent
spectral invariants of individual bundle components. The modulus
$|b| = \sqrt{K/3}$ and the phase $\theta$ are both determined by $K$
alone. The irreducible datum remains $K$ itself.


%======================================================================
\section{What Is and Is Not Established}
%======================================================================

Following the epistemic standards of the GeoVac program, we assess
each link in the hypothetical derivation chain.

\begin{table}[h]
\centering
\begin{tabular}{clc}
\hline\hline
Link & Description & Status \\
\hline
1 & Hopf bundle as geometric arena & $\checkmark$ \\
2 & Spectral data $\to$ $B$, $F$, $\Delta$ & $\bullet$\footnotemark[1] \\
3 & Combination rule $K = \pi(B+F-\Delta)$ & $\times$ \\
4 & $\mathbb{Z}_3$ circulant $\to$ cubic & $\circ$\footnotemark[2] \\
5 & $\alpha$ from cubic root & $\checkmark$ \\
\hline\hline
\end{tabular}
\caption{Derivation chain assessment.
$\checkmark$ = established,
$\bullet$ = largely established (see footnote),
$\circ$ = partially established,
$\times$ = not established.}
\label{tab:chain}
\end{table}
\footnotetext[1]{The selection principle $B/N = 3$ is algebraically
proven (Sec.~III), and the Casimir identity
$\langle C_3\rangle_n = C_4(n)$ connects $B$ to the $SO(4)$
representation theory. The remaining gap: why the degeneracy-weighted
Casimir trace specifically (rather than some other spectral quantity)
for each bundle component. The second selection principle
(Sec.~\ref{sec:second_selection}) provides independent support.}
\footnotetext[2]{The Hermiticity of the circulant is now proven:
$b$ and $c$ are necessarily complex conjugates (Sec.~VI\,D). This is
consistent with $U(1)$ gauge structure but does not yet derive
$bc = K/3$.}

\textbf{Link~1} (established): The Hopf fibration is the unique
principal $U(1)$-bundle over $S^2$ with $c_1 = 1$. That $S^3$ is the
electron state manifold is proven in Ref.~\cite{paper7}. Using the
Hopf bundle as a geometric framework is a natural choice.

\textbf{Link~2} (largely established): The selection principle
$B/N = 3$ is proven algebraically to select $n_{\max} = 3$ uniquely
(Sec.~III). The per-shell Casimir identity
$\langle C_3\rangle_n = C_4(n) = (n^2-1)/2$ establishes that $B$ is
the trace of the $SO(3)$ Casimir over the truncated Peter--Weyl
decomposition of $L^2(S^3)$, connecting it to the representation
theory of $SO(4)$. The remaining gap is the \textit{choice} of which
spectral quantity to extract from each bundle component: why the
Casimir trace for the base (not, say, the spectral determinant), why
$\zeta(2)$ for the fiber (not $\zeta(1)$ or $\zeta(3)$), and why
$1/(|\lambda_3| \times N(2))$ for the boundary.

\textbf{Link~3} (not established): This is the critical gap. No
principle selects the specific additive combination
$K = \pi(B + F - \Delta)$ over alternatives such as
$K = \pi B + F - \Delta$ or $K = \pi B F / \Delta$. The overall
prefactor $\pi$ could be argued as natural (it appears in the
stereographic projection), but the additive structure is unexplained.
A systematic search over spectral determinants, Ray--Singer analytic
torsion, Seeley--DeWitt heat kernel coefficients, and
Atiyah--Patodi--Singer invariants of the three bundle components finds
no natural combination that reproduces $K$ or $K/\pi$. The spectral
determinant product
$\det'(\Delta_{S^1}) \cdot \det'(\Delta_{S^3})/\pi \approx 41.957$
approaches $B = 42$ to $0.1\%$ (see Sec.~VIII) but the gap is not
bridgeable by known spectral corrections. The combination rule mixes a
discrete combinatorial quantity ($B = 42$) with a continuous spectral
invariant ($F = \zeta(2)$) in a way that has no precedent in spectral
geometry. If Link~3 can be closed, the mechanism likely involves the
representation theory of $SO(4)$ rather than the analysis of smooth
differential operators.

\textbf{Link~4} (partial): The cubic form arises naturally from
$\mathbb{Z}_3$-symmetric coupling, and the $SU(2)$ structure constants
share this symmetry. But no derivation connects the Hopf bundle to the
specific circulant with $bc = K/3$.

\textbf{Link~5} (established): Given the cubic, selecting the smallest
positive root is pure algebra.

The structural choices that remain underived are:
\begin{enumerate}
\item Why additive combination $B + F - \Delta$ (not multiplicative)?
\item Why $\pi$ as the overall prefactor?
\item Why $\zeta(2)$ specifically for the fiber (not $\zeta(1)$,
  $\zeta(3)$, etc.)?
\item Why $1/(|\lambda_3| \times N(2))$ for the boundary (not other
  spectral ratios)?
\item Why $\det(M) = -1$?
\end{enumerate}

\textit{Three-tier composition.} $K = \pi(B + F - \Delta)$ assembles
three categorically different objects: $B$ (finite Casimir, rational),
$F$ (infinite Dirichlet series at the packing exponent,
transcendental), $\Delta$ (boundary product at the truncation edge,
rational). Computation identifies each piece but finds no common
generator. The open question is reframed from ``how is each piece
derived'' to ``why does their additive combination equal
$\alpha^{-1}$ to $8.8\times 10^{-8}$''---a structural conspiracy
about the sum, not a derivation about the parts.


%======================================================================
\section{Discussion}
%======================================================================

\subsection{Comparison with Wyler's formula}

Wyler's 1969 formula~\cite{wyler1969} also invokes symmetric spaces
and achieves $\sim 10^{-5}$ precision. However, Wyler's derivation
was criticized for lacking a physical mechanism connecting the
symmetric space volumes to electromagnetic
coupling~\cite{robertson1971}. The present work improves on this
situation in two respects: (i)~$S^3$ is the \textit{proven} electron
state manifold, providing a physical mechanism, and (ii)~the
$p$-value of $5.2 \times 10^{-9}$ quantitatively distinguishes the
formula from numerology. However, both approaches share the weakness
of an unjustified combination rule.

\subsection{Comparison with previous approach}

An earlier attempt by the same author~\cite{paper2_old} modeled
$\alpha^{-1}$ as a ``symplectic impedance'' ratio on the
paraboloid lattice, achieving $6.2 \times 10^{-5}$ relative error
with a fitted helical pitch. The present result achieves
$8.8 \times 10^{-8}$---a $500\times$ improvement---with zero free
parameters.

\begin{table}[h]
\centering
\begin{tabular}{lccc}
\hline\hline
Method & Error & Free params & Cutoff \\
\hline
Symplectic impedance~\cite{paper2_old} & $6.2\times 10^{-5}$ & 1 & $n=5$ \\
\textbf{Hopf bundle (this work)} & $\mathbf{8.8\times 10^{-8}}$ & \textbf{0} & $n=3$ \\
\hline\hline
\end{tabular}
\caption{Evolution of the geometric approach.}
\label{tab:evolution}
\end{table}


\subsection{Spectral determinant near-miss}

The zeta-regularized spectral determinants of the bundle components
are $\det'(\Delta_{S^1}) = 4\pi^2$,
$\det'(\Delta_{S^2}) = e^{1/2 - 4\zeta'_R(-1)}$, and
$\det'(\Delta_{S^3}) = \pi\, e^{\zeta(3)/(2\pi^2)}$. The product
\begin{equation}
  \frac{\det'(\Delta_{S^1}) \cdot \det'(\Delta_{S^3})}{\pi}
  = 4\pi^2 e^{\zeta(3)/(2\pi^2)} \approx 41.957
\end{equation}
is $0.1\%$ from $B = 42$. Whether this near-miss reflects a deeper
connection between the Casimir trace and spectral determinants, or is
merely a coincidence, remains an open question. If a correction term
derivable from the fibration geometry could bridge the $0.1\%$ gap
between $4\pi^2 e^{\zeta(3)/(2\pi^2)}$ and~$42$, it would provide a
path from standard spectral determinants to the Casimir trace,
potentially grounding Link~2 of the derivation chain
(Table~\ref{tab:chain}) in established spectral geometry. Replacing $B = 42$
with this determinant product makes the agreement with CODATA
$11{,}000\times$ \textit{worse}, confirming that the integer Casimir
trace is the correct object, not the spectral determinant.

A systematic search over additive corrections to the spectral
determinant product finds the best candidates to be $+\pi/72$ (0.002\%
residual) and $+1/24$ (0.003\% residual), but neither is derivable from
first principles~\cite{phase2_kappa}. The quantity $\pi/72 =
\pi/(|\lambda_3| \cdot n_{\max}^2)$ could be a finite-size correction,
and $1/24$ appears in heat kernel coefficients, but both remain
\textit{ad hoc}. A broader search over simple arithmetic combinations of
spectral determinants and fibration data finds no exact bridge. This
strengthens the conclusion that the integer Casimir trace $B = 42$---not
the spectral determinant---is the correct object for the base invariant.


\subsection{Algebraic structure of $\alpha$}

The formula implies that $\alpha$ is algebraic over the field
$\mathbb{Q}(\pi)$, since $\zeta(2) = \pi^2/6$. That is,
$\alpha$ is a root of a polynomial whose coefficients are rational
functions of $\pi$ alone. If confirmed, this would constrain $\alpha$
far more tightly than the Standard Model currently does, and would
imply that $\alpha$ is transcendental (since $\pi$ is).


\subsection{$S^3$ specificity}\label{sec:s3_specificity}

\textit{A universal algebraic identity.}
For any round sphere $S^d$, define the \textit{formal} spectral moment
\begin{equation}\label{eq:B_formal}
  B_{\mathrm{formal}}(m) \;=\;
  \tfrac{1}{2}\sum_{k=1}^{m} g_k\,|\lambda_k|,
  \qquad
  N(m) \;=\; \sum_{k=1}^{m} g_k,
\end{equation}
where $\lambda_k = -k(k+d-1)$ and
$g_k = \binom{k+d}{d} - \binom{k+d-2}{d}$ are the eigenvalues and
degeneracies of the Laplace--Beltrami operator on $S^d$. At $m = 2$
these evaluate to $B_{\mathrm{formal}}(2) = \tfrac{1}{2}(d+1)d(d+4)$
and $N(2) = (d+1)(d+4)/2$, giving
\begin{equation}\label{eq:universal_identity}
  \frac{B_{\mathrm{formal}}(2)}{N(2)} \;=\; d
  \qquad\text{for all } d \geq 1.
\end{equation}
Setting $B_{\mathrm{formal}}(m)/N(m) = d$ yields the quadratic
$m^2 + dm - 2(d+2) = 0$, whose unique positive root is $m = 2$
regardless of~$d$. The dimension-matching condition is therefore a
universal property of eigenvalue moments on round spheres, not a
selection principle specific to~$S^3$.

\textit{Why the formula is specific to $S^3$.}
The quantity $B$ appearing in the formula (Sec.~\ref{sec:selection})
is \textit{not} $B_{\mathrm{formal}}$. It is the
degeneracy-weighted trace of the $SO(3)$ \textit{sub}-Casimir
$l(l+1)$ over the $SO(4) \to SO(3)$ branching of each $S^3$ shell:
\begin{equation}\label{eq:B_branch}
  B \;=\; \sum_{n=1}^{m}\sum_{l=0}^{n-1}(2l+1)\,l(l+1).
\end{equation}
For $S^3$, $B = B_{\mathrm{formal}}$ because the diagonal
$SO(3) \subset SO(4)$ embedding has index~1, giving the identity
$\langle C_{SO(3)}\rangle_n = n(n+2)/2 = |\lambda_n|/2$. For all
other spheres $S^d$ with $d \neq 3$, $B \neq B_{\mathrm{formal}}$
and $B/N = d$ has no solution at any integer~$m$.

Three independent properties of $S^3$ underlie this coincidence:
\begin{enumerate}
\item $SO(4) \cong (SU(2) \times SU(2))/\mathbb{Z}_2$ is a product
  group, yielding the clean Clebsch--Gordan branching
  $l = 0, \ldots, n{-}1$ within each shell. No analogous product
  structure exists for $SO(d{+}1)$ with $d \neq 3$.
\item $S^3 \cong SU(2)$ is the only sphere that is a Lie group,
  enabling the Peter--Weyl decomposition that makes $B$ computable
  as a representation-theoretic trace.
\item The embedding index of diagonal $SO(3) \subset SO(4)$ is~1.
  For the quaternionic Hopf fibration, the corresponding
  $SO(5) \subset SO(8)$ embedding through $Sp(2) \times Sp(1)$ has
  index ratio $5/14$, which would give $B/N = 5 \neq 7$ if used as
  the branching factor.
\end{enumerate}
The formula is not one instance of a family parametrized by $d$; it
exists only for~$S^3$.

\textit{Quaternionic and octonionic Hopf fibrations.}
Applying the combination rule $K = \pi(B + F - \Delta)$ formally to
all four Hopf fibrations
($S^0 \to S^1 \to S^1$, $S^1 \to S^3 \to S^2$,
$S^3 \to S^7 \to S^4$, $S^7 \to S^{15} \to S^8$) yields
$1/\alpha \approx 15.4$, $137.036$, $970.0$, and $7164.1$
respectively~\cite{phase6_hopf}. Only the complex Hopf value matches
a known coupling constant. The quaternionic value $K/\pi \approx 309$
satisfies $K_{\mathrm{quat}}/K_{\mathrm{em}} \approx 7.08 \approx
\dim(S^7)$, but the $1.1\%$ deviation from~7 and the absence of any
coupling at $\alpha_W \approx 1/29.5$ rule out a weak-interaction
interpretation. The representation-theoretic obstructions above
explain why: without the $SO(4)$ product structure, the
branching-based $B$ that makes the formula work has no natural
generalization.

\textit{Dynamical specificity: the Fock projection rigidity theorem.}
The representation-theoretic arguments above show that the
\textit{algebraic structure} of $B = 42$ is unique to $S^3$. An
independent \textit{dynamical} argument reaches the same conclusion
from the Fock projection itself, rather than from the $SO(4)$
branching combinatorics. The rigidity theorem (established in the
nuclear extension of the GeoVac framework) states that the
stereographic projection $p_0 = \sqrt{-2 E_n}$ maps a one-electron
central-field Hamiltonian onto the free Laplacian on $S^3$ if and
only if the spectrum $E_{nl}$ is independent of $l$ within each
shell~$n$. The Coulomb potential $-Z/r$ is the unique central
potential exhibiting this accidental degeneracy, by the $SO(4)$
symmetry generated by the Runge--Lenz vector. Any deformation
$V(r) = -Z/r + \delta V(r)$ with $\delta V \not\equiv 0$ lifts the
$l$-degeneracy and breaks the conformal equivalence between the
radial Schr\"odinger problem and the $S^3$ graph Laplacian.
The harmonic oscillator, Woods--Saxon, square well, and Yukawa
potentials all share the angular momentum algebra ($SO(3)$ Gaunt
selection rules) of the Coulomb problem, but none of them carries
the Fock projection because none has the $SO(4)$ accidental
degeneracy. The angular sparsity of the GeoVac
framework~\cite{paper22_sparsity} is universal across spherical
fermion systems; the $S^3$ conformal projection and the Hopf bundle
structure built upon it are Coulomb-specific.

The implication for the present conjecture is that any derivation of
$\alpha$ from spectral geometry must use the Coulomb-specific
features of $S^3$ --- the accidental degeneracy, the Runge--Lenz
vector, and the conformal weight assignments that the Fock projection
imposes --- not merely the $SO(3)$ angular momentum structure that
$S^3$ shares with other manifolds. As a numerical check on this
narrowing, the analog of $B = 42$ in the nuclear shell model
($j$-coupled $(2j+1)\,j(j+1)$ trace at the same truncation
$n_{\max} = 3$) evaluates to $103.5$ and shows no coincidence with
any quantity in $K = \pi(B + F - \Delta)$ at the $p < 10^{-3}$
level~\cite{track_nj_alpha_memo}. This is consistent with the
structural argument: nuclear shell model angular momentum lives on
$SO(3)$ with broken $SO(4)$, lacks the Fock projection, and
therefore cannot be expected to produce the same combinatorial
constants as the Coulomb-specific calculation.

The Bargmann--Segal lattice for the harmonic oscillator (Paper~24,
\cite{paper24_bargmann}) sharpens this further: the harmonic
oscillator's natural discretization on the Hardy space of $S^5$ is
demonstrably $\pi$-free in exact rational arithmetic, with a linear
projection $E_N = \hbar\omega(N + 3/2)$ that requires no calibration
exchange constant. The harmonic oscillator's complex / first-order
geometry produces no spectral zeta corrections at all, so no
$\pi$-class transcendental enters its graph-to-physics map. The
appearance of $\pi$ in $K = \pi(B + F - \Delta)$ is therefore not a
generic feature of quantum discretization but a specific consequence
of the Coulomb potential's $S^3$ Riemannian (second-order) character;
the harmonic oscillator's $S^5$ complex (first-order) character has
no analog. Any future derivation of the combination rule must rely on
the Riemannian/$SO(4)$ structure that distinguishes Coulomb from
every other central potential.


\subsection{Transition operators and the Hopf geometry}

The discrete graph Laplacian for hydrogen has two types of lattice
transitions~\cite{phase2_kappa}: $T_\pm$ (changing $n$ by $\pm 1$,
preserving $l$ and $m$) and $L_\pm$ (changing $m$ by $\pm 1$, preserving
$n$ and $l$). We observe that these correspond naturally to the Hopf
fibration:
\begin{itemize}
\item $L_\pm$ is motion along the Hopf fiber $S^1$. The magnetic quantum
  number $m$ is the $U(1)$ charge (eigenvalue of $J_z$); incrementing
  $m$ rotates around the fiber circle.
\item $T_\pm$ is radial motion on $S^3$, perpendicular to the fiber.
  Changing $n$ moves between eigenspaces of $\Delta_{S^3}$; in the Fock
  stereographic picture, this changes the ``latitude'' on~$S^3$.
\end{itemize}
The angular momentum $l$ is conserved by both transitions---it labels
the $SO(3)$ representation within each $SO(4)$ shell but does not
contribute an independent transition type. The base $S^2$ constrains
state labeling (which $l$, $m$ values exist) but does not generate
lattice transitions. The Fock stereographic projection collapses the
2-dimensional base directions into a single radial variable (the
principal quantum number $n$), explaining why $d_{\max} = 2 \times 2 = 4$
(two motion types $\times$ two directions) rather than~6.


\subsection{Second selection principle}\label{sec:second_selection}

The topological maximum degree $d_{\max} = 4$ gives an asymptotic
bipartite spectral radius bound of $2 d_{\max} = 8$ for the graph
Laplacian~\cite{phase1_kappa}. The $S^3$ Laplace--Beltrami eigenvalue
$|\lambda_n| = n^2 - 1$ equals this bound uniquely at $n = 3$:
\begin{equation}\label{eq:second_selection}
  n^2 - 1 = 2\,d_{\max} \quad\Longrightarrow\quad n_{\max} = 3.
\end{equation}
This provides an independent selection of $n_{\max} = 3$, from the
discrete--continuous compatibility condition rather than from the Casimir
selection $B/N = \dim(S^3)$.

We note a subtlety: $d_{\max} = 4$ is the topological maximum degree for
interior nodes of the lattice class (first realized at $n_{\max} = 4$),
not the realized maximum at $n_{\max} = 3$ (which is~$3$, since all 14
nodes are boundary nodes)~\cite{phase2_kappa}. The selection principle
should be read as: the $S^3$ eigenvalue at $n_{\max}$ matches the
\textit{asymptotic} spectral radius bound of the lattice topology.

The connection to $\kappa$ is immediate: if
$\kappa = -E_1/|\lambda_{n_{\max}}| = -1/(2(n^2-1))$, then the Hopf
selection $n_{\max} = 3$ gives $\kappa = -1/16$, the kinetic scale
factor of the GeoVac program~\cite{paper7}.

\textit{Status:} this is a structural observation, not a derivation. The
two selection principles ($B/N = 3$ and $2d_{\max} = n^2 - 1$) are
algebraically independent; whether they share a common origin in the
Hopf fibration geometry is an open question.


\subsection{Open questions}

\begin{enumerate}
\item \textit{Derive the combination rule} from a path integral or
  spectral determinant identity on the Hopf bundle. This is the
  critical gap (Link~3 in Table~\ref{tab:chain}).

\item \textit{Connect the $\mathbb{Z}_3$ circulant} to the
  differential geometry of the Hopf connection. The $SU(2)$ structure
  constants share the cyclic symmetry; can the coupling strength
  $K/3$ be derived from the connection curvature?

\item \textit{Understand the residual $8.8 \times 10^{-8}$.} Higher-order
  spectral invariants (e.g., the $\eta$-invariant of $S^3$, which
  vanishes by symmetry), gravitational corrections, or weak-mixing
  contributions could account for the discrepancy. One-loop QED
  running of the formula's $\alpha$ from the electron mass to $1$~TeV
  shows that the formula's value corresponds to the Thomson limit
  ($q^2 \to 0$), matching CODATA most closely at $\mu \approx m_e$.
  The residual is intrinsic to the formula, not explained by
  renormalization group running.

  An internal structural check~\cite{alpha_sprint_a_memo} finds that the
  post-cubic residual
  $R_{\mathrm{predict}} = K - \alpha^{-1} - \alpha^{2}$ (where the
  $\alpha^{2}$ term is the trivial offset from the cubic) matches
  $\pi^{3} \alpha^{3}$ to within $0.25\%$ at 80-digit precision:
  $R_{\mathrm{predict}} = 1.2079 \times 10^{-5}$,
  $\pi^{3}\alpha^{3} = 1.2049 \times 10^{-5}$,
  $R_{\mathrm{predict}}/(\pi^{3}\alpha^{3}) = 1.0025$. The $\pi^{3}$
  prefactor is uniquely picked out among $\{\pi^{2}, \pi^{3}, \pi^{4}\}$
  (the other candidates are off by factors of $\sim 3$). A natural
  structural reading is $\pi^{3} = \mathrm{Vol}(S^{5})$, connecting to
  the Bargmann--Segal $S^{5}$ lattice of Paper~24~\cite{paper24_bargmann}.
  An internal cross-check computes the Seeley--DeWitt coefficients
  $a_{0}, a_{2}, a_{4}$ on round $S^{5}$ (each factorizes as
  $\pi^{3} \times \text{rational}$) and finds that no such coefficient
  produces a contribution at order $\alpha^{3}$ structurally against
  the leading $\alpha^{-1}$ of $K$: the Connes--Chamseddine asymptotic
  expansion on $d = 5$ is in odd powers of $\Lambda$ only, with no log
  term, and standard spectral-action truncations fix $\alpha$ once via
  $a_{0}/a_{4}$ ratios, not at order $\alpha^{3}$. The
  $\pi^{3}\alpha^{3}$ match is therefore a numerical pattern of
  unexplained structural origin at this precision---a suggestive hint,
  not a derivation. Paper~24's $\pi$-free certificate on the discrete
  Bargmann--Segal graph is reconciled: it applies to the discrete
  spectrum and edge weights, not to continuum integration measures
  such as $\mathrm{Vol}(S^{5})$.

\item \textit{The other cubic roots.} The two non-physical roots of
  Eq.~(\ref{eq:cubic}) satisfy $\alpha \approx \pm 11.7$,
  corresponding to $1/\alpha \approx \pm 0.085$. A systematic
  comparison of all natural transformations of these roots (including
  $r$, $1/r$, $r^2$, $\sin^2(\arctan r)$, etc.) against known
  Standard Model parameters ($\sin^2\theta_W$, $\alpha_s$, Cabibbo
  angle, Yukawa couplings, mass ratios) finds no match at the $1\%$
  level. The other roots do not appear to encode additional coupling
  constants.

\item \textit{Common origin of the two selection principles.} The
  second selection principle ($2d_{\max} = n^2 - 1$ at $n = 3$) and
  the Casimir selection ($B/N = 3$ at $n = 3$) are algebraically
  independent. Whether they share a common origin in the Hopf
  fibration geometry is open.

\item \textit{Circulant phase and the Hopf connection.} The
  Hermiticity of the circulant coupling and its $\theta \approx
  3\pi/2$ phase at the physical root may encode the $U(1)$ gauge
  structure. A derivation connecting the circulant phase to the Hopf
  connection 1-form would ground Link~4 of
  Table~\ref{tab:chain}.
\end{enumerate}


%======================================================================
\section{Conclusion}
%======================================================================

The fine structure constant satisfies the cubic equation
$\alpha^3 - K\alpha + 1 = 0$, where
$K = \pi(42 + \pi^2/6 - 1/40)$ is constructed from three spectral
invariants of the Hopf fibration. The selection principle
$B/N = \dim(S^3)$ is proven algebraically to select $n_{\max} = 3$
uniquely, via the closed-form ratio $B/N = 3(m+2)(m-1)/10$; this
result is specific to~$S^3$: a formal spectral moment
$B_{\mathrm{formal}}/N = d$ holds universally at $m = 2$ for all
round spheres (Sec.~\ref{sec:s3_specificity}), but the
branching-based $B$ that enters the formula coincides with
$B_{\mathrm{formal}}$ only for $d = 3$, due to the unit embedding
index of $SO(3) \subset SO(4)$. The quaternionic and octonionic Hopf
fibrations yield $1/\alpha \approx 970$ and $7164$ respectively,
matching no known coupling constants. The formula contains no fitted parameters and achieves
$8.8 \times 10^{-8}$ relative agreement with experiment. A
combinatorial search over $1.92 \times 10^9$ formulas yields a
$p$-value of $5.2 \times 10^{-9}$, establishing that this precision
is very unlikely to be accidental. The cubic form is the
characteristic polynomial of a traceless $\mathbb{Z}_3$-symmetric
circulant matrix, providing structural justification for the
self-consistency equation.

The circulant coupling is necessarily Hermitian ($b$ and $c$ are
complex conjugates), with the physical root at phase
$\theta \approx 3\pi/2$, consistent with $U(1)$ gauge structure. A
second selection principle---the bipartite spectral radius bound
$2d_{\max} = 8$ matching $|\lambda_3| = n^2 - 1 = 8$---independently
selects $n_{\max} = 3$ and connects to the kinetic scale factor
$\kappa = -1/16$.

The combination rule $K = \pi(B + F - \Delta)$ remains conjectural.
A systematic investigation of spectral determinants, analytic torsion,
heat kernels, and index-theoretic invariants finds no smooth
spectral-geometric derivation; the rule mixes discrete
representation-theoretic data ($B = 42$) with continuous spectral
invariants ($F = \zeta(2)$) in a way that has no known precedent.
The formula has the character of a remarkable empirical observation
with structural support---the spectral ingredients are natural, the
cubic form is natural, the selection principle is proven, but their
assembly into $K$ awaits a first-principles derivation.


\begin{acknowledgments}
This work supersedes an earlier ``symplectic impedance''
approach~\cite{paper2_old} by the same author. The evolution from
fitted to parameter-free is documented for intellectual transparency.
The four-phase computational audit (arithmetic verification,
combinatorial $p$-value analysis, cubic structure investigation, and
derivation chain assessment) that validated the formula and identified
the circulant structure is available as supplementary
material~\cite{phase2_audit}. A systematic investigation of spectral
determinants, analytic torsion, heat kernel expansions, and
index-theoretic invariants of the Hopf bundle components, documented
in the project repository, established that the combination rule
(Link~3) lies outside the scope of standard smooth spectral geometry.
A subsequent three-phase investigation of the kinetic scale factor
$\kappa = -1/16$ established the transition operator--Hopf
correspondence, a second selection principle from the bipartite
spectral radius bound, and the forced Hermiticity of the circulant
coupling~\cite{phase1_kappa,phase2_kappa,phase3_audit}.
A further investigation applied the formula's recipe to all four Hopf
fibrations, discovered the universal algebraic identity
$B_{\mathrm{formal}}/N = d$, and confirmed the $S^3$ specificity of
the branching-based selection
principle~\cite{phase6_hopf}.
Computational implementation and documentation drafting were performed
with the assistance of large language models (Anthropic Claude) under
the author's scientific direction. All results were validated against
analytical benchmarks.
\end{acknowledgments}

\begin{thebibliography}{20}

\bibitem{feynman1985}
R.~P. Feynman, \textit{QED: The Strange Theory of Light and Matter}
(Princeton University Press, 1985).

\bibitem{eddington1935}
A.~S. Eddington, \textit{New Pathways in Science}
(Cambridge University Press, 1935).

\bibitem{barrow2002}
J.~D. Barrow, \textit{The Constants of Nature}
(Vintage Books, 2002).

\bibitem{paper7}
J.~Loutey, ``The Dimensionless Vacuum: $S^3$ Conformal Geometry and
the Schr\"odinger Equation,'' GeoVac Paper~7 (2026).

\bibitem{paper22_sparsity}
J.~Loutey, ``Potential-Independent Angular Sparsity in Spherical
Fermion Systems,'' GeoVac Paper~22 (2026).

\bibitem{paper24_bargmann}
J.~Loutey, ``The Bargmann--Segal Lattice: $\pi$-Free Discretization
of the Harmonic Oscillator on $S^5$,'' GeoVac Paper~24 (2026).

\bibitem{track_nj_alpha_memo}
J.~Loutey, ``Track NJ: Fine Structure Constant and the Nuclear
Extension --- Theory Memo,'' GeoVac internal memo (2026). Available
at \texttt{docs/track\_nj\_alpha\_memo.md} in the project repository.

\bibitem{hopf1931}
H.~Hopf, ``\"Uber die Abbildungen der dreidimensionalen Sph\"are auf
die Kugelfl\"ache,'' \textit{Math.\ Ann.}\ \textbf{104}, 637 (1931).

\bibitem{urbantke2003}
H.~K. Urbantke, ``The Hopf fibration---seven times in physics,''
\textit{J.\ Geom.\ Phys.}\ \textbf{46}, 125 (2003).

\bibitem{fock1935}
V.~Fock, ``Zur Theorie des Wasserstoffatoms,''
\textit{Z.\ Phys.}\ \textbf{98}, 145 (1935).

\bibitem{paper2_old}
J.~Loutey, ``The Fine Structure Constant as Geometric Impedance: A
Symplectic Framework,'' GeoVac Paper~2, v1 (2026). [Superseded by
the present work.]

\bibitem{wyler1969}
A.~Wyler, ``L'espace sym\'etrique du groupe des \'equations de
Maxwell,'' \textit{C.\ R.\ Acad.\ Sci.\ Paris}\ \textbf{269A}, 743
(1969).

\bibitem{robertson1971}
B.~Robertson, ``Wyler's expression for the fine-structure constant,''
\textit{Phys.\ Rev.\ Lett.}\ \textbf{27}, 1545 (1971).

\bibitem{phase2_audit}
J.~Loutey, ``Paper~2 Computational Audit: Phases 1--4,'' GeoVac
supplementary material (2026). Available at
\texttt{debug/alpha\_audit/} in the project repository.

\bibitem{phase1_kappa}
J.~Loutey, ``$\kappa = -1/16$ Investigation: Phase~1,'' GeoVac
supplementary material (2026). Available at
\texttt{debug/kappa\_investigation/} in the project repository.

\bibitem{phase2_kappa}
J.~Loutey, ``Spectral Determinants \& Second Selection Principle:
Phase~2,'' GeoVac supplementary material (2026). Available at
\texttt{debug/kappa\_investigation/phase2/} in the project repository.

\bibitem{phase3_audit}
J.~Loutey, ``Circulant Structure of the Cubic: Phase~3,'' GeoVac
supplementary material (2026). Available at
\texttt{debug/kappa\_investigation/phase3/} in the project repository.

\bibitem{phase6_hopf}
J.~Loutey, ``Quaternionic Hopf Fibration \& Universal Algebraic
Identity: Phase~6,'' GeoVac supplementary material (2026). Available
at \texttt{debug/kappa\_investigation/phase6/} in the project
repository.

\bibitem{marcolli_vs_2014}
M.~Marcolli and W.~D.~van Suijlekom, ``Gauge networks in noncommutative
geometry,'' \textit{J.\ Geom.\ Phys.}\ \textbf{75}, 71 (2014),
arXiv:1301.3480.

\bibitem{perez_sanchez_2024}
C.~I.~Perez-S\'anchez, preprint (2024), arXiv:2401.03705. (Revisits
the continuum limit of Marcolli--van Suijlekom gauge networks;
establishes that the limit is Yang--Mills without Higgs.)

\bibitem{perez_sanchez_2024_comment}
C.~I.~Perez-S\'anchez, Comment, preprint (2025), arXiv:2508.17338.

\bibitem{alpha_sprint_a_memo}
J.~Loutey, ``Paper~2 Sprint~A: Structural consistency verification,''
GeoVac supplementary material (2026). Available at
\texttt{debug/alpha\_sprint\_a/} in the project repository.

\end{thebibliography}

\end{document}
